\hypertarget{python-3.11-faster-cpython-specializing-interpreter-and-adaptive-bytecode}{%
\chapter{Python 3.11: Faster CPython Specializing Interpreter and
Adaptive
Bytecode}\label{python-3.11-faster-cpython-specializing-interpreter-and-adaptive-bytecode}}

\hypertarget{pep-659-specialized-adaptive-interpreter-architecture}{%
\subsection{17.1 PEP 659 Specialized Adaptive Interpreter
Architecture}\label{pep-659-specialized-adaptive-interpreter-architecture}}

Prior to Python 3.11, the bytecode interpreter loop executed
instructions generically. For instance, a
\passthrough{\lstinline!LOAD\_ATTR!} instruction looked up attributes
using hash table lookups every time. PEP 659 introduced a
\textbf{specializing adaptive interpreter} to dramatically reduce lookup
overhead.

\hypertarget{monomorphism-vs.-polymorphism-in-dynamic-execution}{%
\subsubsection{1. Monomorphism vs.~Polymorphism in Dynamic
Execution}\label{monomorphism-vs.-polymorphism-in-dynamic-execution}}

While Python is a dynamic language, execution paths in real-world
application code are highly \textbf{monomorphic} (a specific call site
or attribute access typically processes objects of the exact same type
repeatedly). Generic bytecode instructions (like
\passthrough{\lstinline!LOAD\_ATTR!}) are designed to handle any type of
object, performing expensive CPython namespace dictionary lookups and
method searches on every execution.

\hypertarget{dynamic-instruction-patching}{%
\subsubsection{2. Dynamic Instruction
Patching}\label{dynamic-instruction-patching}}

PEP 659 optimizes hot, monomorphic paths by dynamically modifying the
instruction stream in memory at runtime: 1. \textbf{Generic State}: The
compiler generates generic opcodes
(e.g.~\passthrough{\lstinline!LOAD\_ATTR!}). 2. \textbf{Warm-up \&
Monitoring}: When executed, the generic opcode transitions into an
\textbf{adaptive state} (e.g.,
\passthrough{\lstinline!LOAD\_ATTR\_ADAPTIVE!}). The adaptive
instruction maintains an execution counter initialized to
\passthrough{\lstinline!8!}. 3. \textbf{Specialization}: On every
execution, the interpreter monitors the operand types and decrements the
counter. If the type is consistent when the counter reaches
\passthrough{\lstinline!0!}, the interpreter patches the bytecode in
memory, replacing the adaptive opcode with a \textbf{specialized opcode}
(e.g., \passthrough{\lstinline!LOAD\_ATTR\_INSTANCE\_VALUE!}). 4.
\textbf{Deoptimization}: If the operand type changes at a specialized
call site, the opcode execution fails its fast path validation check. It
branches to a deoptimization routine, restores the generic or adaptive
state, and executes the slow lookup.

\begin{lstlisting}
CPython 3.11 Instruction Specialization Lifecycle:

  [ Generic Opcode ] (LOAD_ATTR)
         |
         | (Executed first time)
         v
  [ Adaptive Opcode ] (LOAD_ATTR_ADAPTIVE, Counter = 8)
         |
         | (Consistent types, Counter reaches 0)
         v
  [ Specialized Opcode ] (LOAD_ATTR_INSTANCE_VALUE)
    /                 \
   / (Type Match)      \ (Type Mismatch)
  v                     v
[ Fast Direct Offset ] [ Deoptimize to Generic / Slow Lookup ]
\end{lstlisting}

\begin{center}\rule{0.5\linewidth}{0.5pt}\end{center}

\hypertarget{specialized-bytecode-opcodes-and-inline-cache-memory-layout}{%
\subsection{17.2 Specialized Bytecode Opcodes and Inline Cache Memory
Layout}\label{specialized-bytecode-opcodes-and-inline-cache-memory-layout}}

CPython reserves extra memory slots directly adjacent to the instruction
bytes inside the code object's instruction array. This contiguous memory
block is called the \textbf{Inline Cache}.

\hypertarget{inline-cache-memory-alignment}{%
\subsubsection{1. Inline Cache Memory
Alignment}\label{inline-cache-memory-alignment}}

Instructions are represented by \passthrough{\lstinline!\_Py\_CODEUNIT!}
elements (16-bit blocks: 8-bit opcode and 8-bit argument). When a
cacheable instruction is compiled, the compiler allocates empty
\passthrough{\lstinline!\_Py\_CODEUNIT!} entries directly following the
instruction to serve as the cache.

\hypertarget{cpython-bytecode-stream-cache-layout}{%
\paragraph{CPython Bytecode Stream Cache
Layout:}\label{cpython-bytecode-stream-cache-layout}}

\begin{lstlisting}
Memory Address:  |  n  |  n+1  |  n+2  |  n+3  |  n+4  |  n+5  |  n+6  |
Bytecode:        [ LOAD_ATTR ] [ Cache Slot 0 ] [ Cache Slot 1 ] [ Cache Slot 2 ]
\end{lstlisting}

When executing \passthrough{\lstinline!LOAD\_ATTR!} at address
\passthrough{\lstinline!n!}, the VM knows that the next 3 code units are
reserved for the inline cache and skips them to decode the next
instruction at address \passthrough{\lstinline!n+4!}.

\hypertarget{c-level-cache-structures-inside-pycore_code.h}{%
\subsubsection{\texorpdfstring{2. C-level Cache Structures inside
\texttt{pycore\_code.h}}{2. C-level Cache Structures inside pycore\_code.h}}\label{c-level-cache-structures-inside-pycore_code.h}}

For attribute lookups (\passthrough{\lstinline!LOAD\_ATTR!}), CPython
maps the inline cache slots to the
\passthrough{\lstinline!\_PyAttrCache!} struct:

\begin{lstlisting}[language=C]
typedef struct {
    uint16_t counter;           /* Adaptive execution counter */
    uint32_t type_version;      /* Object type version pointer (tp_version_tag) */
    uint16_t index;             /* Attribute array offset */
} _PyAttrCache;
\end{lstlisting}

For global variable lookups (\passthrough{\lstinline!LOAD\_GLOBAL!}),
the VM allocates the \passthrough{\lstinline!\_PyLoadGlobalCache!}
struct:

\begin{lstlisting}[language=C]
typedef struct {
    uint16_t counter;           /* Adaptive counter */
    uint32_t module_keys_version; /* Dictionary keys version of module globals */
    uint32_t builtin_keys_version; /* Dictionary keys version of builtins dict */
} _PyLoadGlobalCache;
\end{lstlisting}

\hypertarget{cpython-type-versioning-tp_version_tag}{%
\subsubsection{\texorpdfstring{3. CPython Type Versioning
(\texttt{tp\_version\_tag})}{3. CPython Type Versioning (tp\_version\_tag)}}\label{cpython-type-versioning-tp_version_tag}}

To ensure cache validity: * Every class
(\passthrough{\lstinline!PyTypeObject!}) in CPython has an internal
\passthrough{\lstinline!tp\_version\_tag!} field (a 32-bit unsigned
integer). * If a class is modified at runtime (e.g.~adding a method or
setting a class variable), CPython increments a global type version
counter and assigns it to the class's
\passthrough{\lstinline!tp\_version\_tag!}, immediately invalidating all
inline caches containing the old version tag.

\hypertarget{execution-flow-of-specialized-instructions}{%
\subsubsection{4. Execution Flow of Specialized
Instructions}\label{execution-flow-of-specialized-instructions}}

\begin{itemize}
\tightlist
\item
  \textbf{\passthrough{\lstinline!LOAD\_ATTR\_INSTANCE\_VALUE!}}:

  \begin{enumerate}
  \def\labelenumi{\arabic{enumi}.}
  \tightlist
  \item
    Pops the target object off the stack.
  \item
    Compares the object's class
    \passthrough{\lstinline!tp\_version\_tag!} against the cached
    \passthrough{\lstinline!type\_version!}.
  \item
    If they match, it skips dict lookups and reads the attribute
    directly from the object's instance values array
    (\passthrough{\lstinline!obj->obj\_values[index]!}), pushing it onto
    the stack.
  \end{enumerate}
\item
  \textbf{\passthrough{\lstinline!LOAD\_GLOBAL\_MODULE!}}:

  \begin{enumerate}
  \def\labelenumi{\arabic{enumi}.}
  \tightlist
  \item
    Compares the module dictionary's keys version against
    \passthrough{\lstinline!module\_keys\_version!}.
  \item
    If they match, it reads the value directly from the dictionary value
    array, completely bypassing hash collisions and key comparisons.
  \end{enumerate}
\end{itemize}

\begin{center}\rule{0.5\linewidth}{0.5pt}\end{center}

\hypertarget{polymorphic-deoptimization-costs}{%
\subsection{17.3 Polymorphic Deoptimization
Costs}\label{polymorphic-deoptimization-costs}}

While specialization yields massive performance improvements, it is
sensitive to code polymorphism: * \textbf{Deoptimization Paths}: If a
call site receives objects of varying classes (e.g.~calling a function
with both \passthrough{\lstinline!Dog!} and
\passthrough{\lstinline!Cat!} classes where both define
\passthrough{\lstinline!.x!} but at different offsets), the instruction
will constantly deoptimize. * \textbf{Thrashing}: The VM will oscillate
between deoptimizing, executing slow-path lookups, warming up the
adaptive counter, and specializing again (thrashing). This destroys
performance because the interpreter incurs the combined costs of: 1.
Cache tag validation mismatch checks. 2. Generic name dictionary hash
lookups. 3. Adaptive state tracking and bytecode memory rewriting
overhead.

\hypertarget{coding-guidelines-for-pep-659-optimization}{%
\subsubsection{Coding Guidelines for PEP 659
Optimization:}\label{coding-guidelines-for-pep-659-optimization}}

To maximize CPython 3.11+ execution speed, developers should design code
to be \textbf{monomorphic}: 1. \textbf{Keep Types Consistent}: Avoid
passing objects of different types to the same hot functions or
attribute access sites. 2. \textbf{Favor Class Structure Consistency}:
Avoid dynamically modifying instance attributes or class namespaces at
runtime, as this increments \passthrough{\lstinline!tp\_version\_tag!}
and invalidates caches globally.

\begin{center}\rule{0.5\linewidth}{0.5pt}\end{center}
